\section{The first heat correction and mass at high zero indices}
\label{sec:analytic-heat-correction}

\begin{theorem}[NN-A019: Gaussian saddle and first character correction]
Retain the notation of NN-A015. For sufficiently large $t$, let $b_t>1$
be the unique large solution of
\[
 4\pi e^{4b_t}=2tb_t+9,
 \qquad \lambda_t=8tb_t-2t+36.
\]
Let $U_t$ have density $e^{tu^2}\Phi(u)/H_t(0)$ on $[0,\infty)$,
and put
\[
 \mu_t=\mathbb E U_t^2=2p_1(t),\qquad
 v_t=\frac{\operatorname{Var}(U_t^2)}{\mu_t^2}.
\]
The following assertions about the kernel hold without A or M:
\begin{align}
 b_t&=\tfrac14\bigl(\log t+\log\log t-\log(8\pi)+o(1)\bigr),
 &p_1(t)&=\tfrac12b_t^2+o(1),\\
 \sqrt{\lambda_t}(U_t-b_t)&\Longrightarrow\mathcal N(0,1),
 &v_t&\sim\frac{1}{2tb_t^3},\\
 p_1'(t)&\sim\frac{b_t}{4t}\sim\frac{\log t}{16t}.
 \label{eq:heat-first-correction-scale}
\end{align}
The normal limit holds with every absolute moment and every fixed
exponential weight. Locally uniformly for $w\in\mathbb C$,
\begin{equation}\label{eq:heat-first-correction-entire}
 E_t(w)=E_*(w)+\frac{v_t}{4}w\bigl(E_*(w)-E_*'(w)\bigr)+o(v_t),
 \qquad E_*(w)=\cosh\sqrt{2w}.
\end{equation}
Equivalently, with $\epsilon_t=v_t/4$,
\[
 E_t(w)=e^{\epsilon_tw}E_*((1-\epsilon_t)w)+o(\epsilon_t).
\]

Under A, the positive normalized zero weights satisfy, for every
fixed $j$,
\begin{equation}\label{eq:heat-first-correction-weights}
 \beta_j(t)=\beta_j^*(1-v_t/4)+o(v_t),\qquad
 \beta_j^*=\frac{8}{\pi^2(2j-1)^2}.
\end{equation}
For every fixed $k\ge2$, their power sums obey
\[
 \sum_j\beta_j(t)^k
 =\left(1-\frac{k v_t}{4}\right)\sum_j(\beta_j^*)^k+o(v_t).
\]
If $\sigma$ is a fixed finite permutation moving $d(\sigma)$ points,
then its Thoma character has the sharp relative expansion
\begin{equation}\label{eq:heat-first-correction-character}
 \chi_{S,t}(\sigma)=\chi_{S,*}(\sigma)
           \left(1-\frac{d(\sigma)v_t}{4}+o(v_t)\right).
\end{equation}
Nevertheless, the vectors $(\beta(t)-\beta^*)/v_t$ do not converge
in $\ell^1$, and
\[
 \liminf_{t\to\infty}
       \frac{\|\beta(t)-\beta^*\|_1}{v_t}\ge\frac12.
\]
Thus the first-order finite-character limit cannot be obtained by
differentiating the full weight sequence in $\ell^1$.
\end{theorem}

\begin{proof}
Write
\[
 f_t(u)=tu^2+9u-\pi e^{4u},\qquad
 h(u)=\Phi(u)e^{-9u+\pi e^{4u}}.
\]
The exact kernel and its bounds in NN-A015 give constants
$0<c<C<\infty$ with $c\le h(u)\le C$ for all $u\ge0$ and
$h(u)\to2\pi^2$ as $u\to\infty$. The latter follows by factoring
the $n=1$ summand; its negative correction is
$-3\pi e^{-4u}$ and the other summands tend to zero uniformly by
the same summable upper bound used in (A015.7).

For large $t$, $f_t'(1)>0$ whereas $f_t'(u)\to-\infty$.
Since $(f_t')''(u)=-64\pi e^{4u}<0$, there is exactly one root
$b_t>1$, the local maximum. Evaluating $4\pi e^{4b}-2tb-9$ at
$b=(\log t)/4$ and $b=(\log t+2\log\log t)/4$ shows
$b_t\sim(\log t)/4$. Taking logarithms of its defining equation
then gives the more precise displayed expansion. Also
$f_t''(b_t)=-\lambda_t$ and $\lambda_t\sim8tb_t$.

Here is a full tail justification for the Gaussian saddle method.
Fix a small constant $\delta>0$. For all sufficiently large $t$,
\[
 f_t''(u)\le-c_0\lambda_t\qquad (u\ge b_t-\delta)
\]
with $c_0>0$ independent of $t$: insert
$16\pi e^{4b_t}=8tb_t+36$ and use $b_t\to\infty$.
Since $f_t'(b_t)=0$, this implies
\[
 f_t(u)-f_t(b_t)\le-\tfrac12c_0\lambda_t(u-b_t)^2
               \qquad(u\ge b_t-\delta).
\]
On $[0,b_t-\delta]$, the maximum of $f_t$ is at an endpoint.
Indeed $f_t'$ is strictly concave, negative at $0$, positive at $1$,
and zero at $b_t$; its earlier zero is a local minimum.
Moreover $f_t(b_t)\sim tb_t^2$ whereas $f_t(0)=-\pi$.
Consequently
\[
 \sup_{0\le u\le b_t-\delta}f_t(u)
                       \le f_t(b_t)-c_1\lambda_t
\]
for some $c_1>0$. The integral on this interval remains negligible
even after multiplication by any fixed power of
$|\sqrt{\lambda_t}(u-b_t)|$ and by
$\exp(A|\sqrt{\lambda_t}(u-b_t)|)$, for fixed $A$:
its logarithmic extra cost is at most
$A b_t\sqrt{\lambda_t}+O(\log(b_t\lambda_t))=o(\lambda_t)$.

For bounded $y$, Taylor's formula gives
\[
 f_t(b_t+y/\sqrt{\lambda_t})-f_t(b_t)
                  =-y^2/2+O(|y|^3/\sqrt{\lambda_t}),
\]
because $f_t'''=O(\lambda_t)$ in a fixed neighbourhood of $b_t$.
The bounded positive factor $h$ tends there to $2\pi^2$.
The preceding Gaussian majorant and left-tail estimate justify
dominated convergence after this change of variables, including
every fixed polynomial and exponential weight. Hence, writing
$Y_t=\sqrt{\lambda_t}(U_t-b_t)$, one has
$\mathbb E Y_t\to0$, $\mathbb E Y_t^2\to1$, and uniformly bounded
absolute moments of every fixed order with exponential weights.

It follows that
\[
 \mu_t=b_t^2+o(b_t/\sqrt{\lambda_t})+O(1/\lambda_t)
          =b_t^2+o(1),
 \qquad
 \operatorname{Var}(U_t^2)\sim\frac{4b_t^2}{\lambda_t}.
\]
For the variance, expand
$U_t^2=b_t^2+2b_tY_t/\sqrt{\lambda_t}+Y_t^2/\lambda_t$;
the cross term is $O(b_t/\lambda_t^{3/2})$ and the last variance
is $O(1/\lambda_t^2)$, both smaller than the leading term.
The claims about $v_t$ and $p_1$ follow. Differentiating the
positive exponential tilt gives
$p_1'(t)=\operatorname{Var}(U_t^2)/2$, proving its asymptotic.

Set $V_t=U_t^2/\mu_t$, so that $\mathbb E V_t=1$ and
$\mathbb E(V_t-1)^2=v_t$. For each compact set $|w|\le R$,
the entire function of $v$
$g_w(v)=\cosh\sqrt{2wv}$ has, on $v\ge0$, third derivative bounded
by $C_R\exp(C_R\sqrt v)$. This follows either from its power
series on $0\le v\le1$ or from its explicit derivatives on
$v\ge1$. Taylor expansion at $v=1$ therefore has an expected
remainder bounded by
\[
 C_R\mathbb E\bigl[|V_t-1|^3
                         e^{C_R\sqrt{\max(V_t,1)}}\bigr].
\]
To bound it, put $\eta_t=1/(b_t\sqrt{\lambda_t})$. The exact
formula for $U_t^2$ and $\mu_t\sim b_t^2$ show
\[
 |V_t-1|\le C\eta_t(|Y_t|+\mathbb E|Y_t|)
       +C\eta_t^2(Y_t^2+\mathbb E Y_t^2),
 \qquad
 \sqrt{V_t}\le C+C\eta_t|Y_t|.
\]
The weighted moment bounds imply that the remainder is
$O_R(\eta_t^3)=o(v_t)$, since $v_t\sim4\eta_t^2$.
The linear Taylor term has zero mean. With $a=\sqrt{2w}$,
\[
 g_w''(1)=\frac{a^2\cosh a-a\sinh a}{4}
          =\frac w2\bigl(E_*(w)-E_*'(w)\bigr).
\]
Using $E_t(w)=\mathbb E g_w(V_t)$ proves
\eqref{eq:heat-first-correction-entire}, uniformly on the compact set.

Under A, the simple limiting roots
$w_j^*=-1/\beta_j^*$ correspond to the ordered finite-time roots
$w_j(t)=-1/\beta_j(t)$, as proved in NN-A015. The local root
expansion from \eqref{eq:heat-first-correction-entire} is
\[
 w_j(t)-w_j^*
 =-\frac{v_t}{4}\,
       \frac{w_j^*(E_*(w_j^*)-E_*'(w_j^*))}{E_*'(w_j^*)}
       +o(v_t)
 =\frac{v_t}{4}w_j^*+o(v_t).
\]
This proves \eqref{eq:heat-first-correction-weights}. On a fixed
small disc avoiding the zeros of $E_*$,
\[
 \log E_t(w)=\log E_*(w)+\frac{v_t}{4}
                w\left(1-\frac{E_*'(w)}{E_*(w)}\right)+o(v_t).
\]
Cauchy's formula permits coefficient comparison. The coefficient
of $w$ remains one; for $k\ge2$, comparison with
$\log E_t=\sum_{k\ge1}(-1)^{k-1}(\sum_j\beta_j(t)^k)w^k/k$
gives the power-sum formula. Multiplying the finitely many nonzero
cycle values gives \eqref{eq:heat-first-correction-character}.

Finally, if $(\beta(t)-\beta^*)/v_t$ converged in $\ell^1$, every
coordinate of its limit would be $-\beta_j^*/4$ by the proved
fixed-coordinate formula, so its coordinate sum would be $-1/4$.
Continuity of summation on $\ell^1$ instead makes that sum zero,
because each finite-time vector has total mass one. This is a
contradiction. For any fixed $J$, the zero total mass of
$\beta(t)-\beta^*$ also gives
\[
 \|\beta(t)-\beta^*\|_1
 \ge2\left|\sum_{j\le J}(\beta_j(t)-\beta_j^*)\right|.
\]
Divide by $v_t$, take a lower limit, and then let $J\to\infty$ to
obtain the lower bound $1/2$.
\end{proof}

\paragraph{Exact representation-theoretic translation.}
For every sufficiently small $\epsilon>0$,
\[
 e^{\epsilon w}E_*((1-\epsilon)w)
 =e^{\epsilon w}\prod_j(1+(1-\epsilon)\beta_j^*w)
\]
is the generating function of the genuine Thoma parameters
$\alpha=0$, $\beta=(1-\epsilon)\beta^*$ and $\gamma=\epsilon$.
Their total mass is one. Formula
\eqref{eq:heat-first-correction-entire} proves that this path gives
the first-order values of every fixed finite character observable
of the actual heat family, although that family has $\gamma=0$
at every finite time. The missing first-order mass at each fixed
zero index is compensated at indices tending to infinity.
The proved failure of an $\ell^1$ derivative makes this distinction
precise. The corresponding genuine $U(\infty)$ character is
\[
 e^{\epsilon\operatorname{Tr}(U-I)}
 \prod_j\det\bigl((1-(1-\epsilon)\beta_j^*)I
                         +(1-\epsilon)\beta_j^*U\bigr).
\]
Its exponential factor is positive definite on $U(d)$ by the
Poisson expansion in normalized tensor powers of the defining
character, and is compatible under $U\mapsto U\oplus1$.
Consequently the entire-function expansion gives, by finite
products, the first-order unitary character values as well.
These are exact infinite-group and analytic statements, not an
integral positive basis or an arithmetic motive. The only imported
classification is the same primary Thoma theorem used in NN-A008;
the saddle and correction proofs above are complete.
For the exterior-power dictionary, the concrete Grothendieck
locator is SGA~6, Expos\'e~V, Example~2.2(b) and Lemma~2.2.1.
The latter gives the exterior-power filtration for a locally
split exact sequence; the representation-ring case is explicitly
included after the lemma. It does not turn these analytic weights
into integral vector-bundle multiplicities.
